\documentclass[letterpaper,10pt,conference]{ieeeconf}

\IEEEoverridecommandlockouts
\overrideIEEEmargins

\usepackage{amsmath,amssymb}
\usepackage{booktabs}
\usepackage{cite}
\usepackage{graphicx}
\usepackage{microtype}
\usepackage{url}

% Switch to \anonymousfalse for the camera-ready version.
\newif\ifanonymous
\anonymoustrue

\newcommand{\system}{\textsc{Scan2VLA}}
\newcommand{\todo}[1]{\textbf{[TODO: #1]}}
\newcommand{\piHalf}{\ensuremath{\pi_{0.5}}}

\title{\system: Scene-Native Vision-Language-Action Learning with
Interactive Gaussian Digital Twins}

\ifanonymous
  \author{Anonymous Authors}
\else
  \author{Author One$^{1}$, Author Two$^{1}$, and Author Three$^{2}$%
    \thanks{$^{1}$Department of Robotics, University A, City, Country.
      {\tt\small \{author.one,author.two\}@example.edu}}%
    \thanks{$^{2}$Institute B, City, Country.
      {\tt\small author.three@example.edu}}%
    \thanks{This work was supported by \todo{funding source}.}%
  }
\fi

\begin{document}

\maketitle
\thispagestyle{empty}
\pagestyle{empty}

\begin{abstract}
Adapting a general-purpose vision-language-action (VLA) foundation model to a
new robot embodiment remains data-intensive because its cameras, proprioception,
and action semantics differ from pretraining. Photorealistic simulation can
supply adaptation trajectories, but composing a reconstructed background with
independently acquired foreground assets introduces appearance and alignment
errors around the objects on which manipulation depends. We present \system{}, an
embodiment-agnostic pipeline that turns a jointly reconstructed 3D Gaussian
Splatting (3DGS) scene into an interactive Gaussian digital twin. \system{} fuses
multi-view instance evidence into Gaussian labels, separately completes newly
exposed background and unobserved object surfaces, and binds each visual asset to
aligned watertight collision geometry. We evaluate two general-purpose VLA
backbones, \piHalf{} and SmolVLA, on an AgiBot G2 that is absent from their
pretraining and prior adaptation data before our study.
Tasks range from placing one beverage bottle into a picnic basket to whole-body
picnic packing and trunk placement. Simulator construction still uses real-scene
capture, robot models, and G2 calibration; ``zero real demonstrations'' refers
only to post-training adaptation data. At zero real demonstrations, \system-only
achieves \todo{$S/N$ G2 successes with a 95\% confidence interval, versus the
strongest matched simulation-only baseline's $S_b/N_b$}; a paired sim--real
subset measures whether simulation performance predicts G2 performance, and at a
prespecified real-world success threshold \system{} reduces required real
demonstrations by \todo{$X\%$ with a 95\% bootstrap confidence interval}.
These results \todo{state the bounded conclusion supported for this unseen
embodiment, the evaluated tasks/scenes, and the two VLA backbones}.
\end{abstract}

% Target allocation under the ICRA 2027 eight-page total limit, including
% references: Abstract + Introduction 1.15; Related Work 0.40; Method
% (including formulation) 2.25; Experiments 3.20; Limitations + Conclusion
% 0.30; References 0.70 pages. Within Experiments: setup/controls 0.45;
% asset evidence 0.40; sim--real correlation 0.40; demonstration efficiency
% 0.65; long-horizon results 0.65; ablations/generalization 0.35; and
% efficiency/safety/statistics 0.30 pages. Move the 54-cell ledger,
% paired-episode manifest, full per-object tables, and extended failures to
% supplement.

\section{Introduction}
\label{sec:introduction}

VLA policies provide a scalable interface from visual observations and language
instructions to robot actions~\cite{black2025pi05,shukor2025smolvla}. Adapting
such a policy to a new embodiment, however, still requires demonstrations that
cover its cameras, proprioceptive state, action semantics, and target tasks.
Collecting this coverage on hardware consumes robot time, human supervision,
physical resets, and safety monitoring; the cost grows further for whole-body,
long-horizon tasks in which one failed stage invalidates the episode.

Photorealistic simulation can replace part of this collection only if its
observations preserve the cues used by the policy. A common real-to-sim pipeline
reconstructs a static deployment scene, separately captures or generates each
interactive object, and inserts that asset into the scene. Differences in scale,
pose, exposure, color response, and baked illumination then concentrate around
the manipulated object and its boundary. This asset-source gap changes the
policy input precisely during grasping and contact, so an otherwise correct
simulated trajectory can provide visually mismatched supervision.

Recent 3DGS simulators improve rendering fidelity, physics integration, and
throughput~\cite{qureshi2025splatsim,jiang2025gsworld,jia2026gsplayground}; other
systems use reconstructed or hybrid-rendered data for VLA post-training
~\cite{gao2026mirage2matter,kim2026legs}. Scene-level methods can also lift
semantics and extract collision geometry from 3DGS~\cite{sun2026digitaltwin}.
These advances establish the value of photorealistic simulation, but leave an
evaluation gap: it remains unclear whether deriving a movable object's visual
representation from the same joint reconstruction as its background makes the
resulting trajectories more valuable for VLA adaptation.

Our key insight is that the target scene should serve as its own interactive
asset inventory. Based on this insight, \system{} reconstructs the workspace
once, fuses multi-view instance evidence into the scene Gaussians, and extracts
each rigid object as a movable subset in the scene coordinate and capture domain.
It separately completes the background exposed by object removal and the object
surfaces exposed by motion, generates aligned watertight collision geometry, and
binds the asset to a physics engine. Figure~\ref{fig:teaser} summarizes this
scan-to-skill pipeline.

\begin{figure*}[t]
  \centering
  \fbox{\parbox[c][3.0cm][c]{0.96\textwidth}{\centering
    Teaser: general-purpose VLA foundation model $\rightarrow$ calibrated G2
    scene capture $\rightarrow$ native interactive Gaussian twin
    $\rightarrow$ simulation-only adaptation $\rightarrow$ first-time execution
    on an unseen AgiBot G2. Include one
    side-by-side comparison with an independently scanned object, plus G2
    rollouts for the bottle-to-basket probe and primary picnic-to-trunk task.}}
  \caption{Given one calibrated capture, \system{} extracts movable objects and
  collision geometry from the same native 3DGS as the background. The matched
  baseline reconstructs the object independently before insertion. Policies are
  post-trained on identical trajectories and evaluated on AgiBot G2, which is
  absent from the foundation models' pretraining data, from a
  bottle-to-basket pick-and-place to picnic packing and trunk placement. The final panel
  reports \todo{the representative Scan2VLA-only and strongest-baseline success
  rates, including trial counts}.}
  \label{fig:teaser}
\end{figure*}

We evaluate this hypothesis on the AgiBot G2 using the general-purpose VLA
foundation models \piHalf{} and SmolVLA. G2 is absent from their pretraining
data, and the starting checkpoints receive no G2-specific policy adaptation data;
the
robot models, calibration, action interface, and low-level controller are
G2-specific but frozen across conditions. The task suite uses a beverage-bottle
pick-and-place into a picnic basket as a diagnostic probe and centers on a
long-horizon task that packs tabletop snack
items into a picnic basket before G2 uses whole-body motion to lift the loaded
basket and place it in a car trunk. At zero real demonstrations, \system{} obtains
\todo{$S/N$ end-to-end successes compared with $S_b/N_b$ for the strongest
matched simulation-only baseline}; along the real-data learning curve, it saves
\todo{$X\%$ of real demonstrations at threshold $\tau$}. These experiments
support three contributions:
\begin{itemize}
  \item An embodiment-agnostic scene-native decomposition method that fuses
  multi-view instance evidence into a jointly optimized 3DGS and extracts rigid
  objects without a separate object scan or cross-asset appearance alignment.
  \item An interactive Gaussian digital-twin construction procedure that
  separately completes the exposed background and unobserved object surfaces,
  preserves provenance for observed and synthesized Gaussians, and produces
  aligned watertight collision geometry. Physical-property estimation remains a
  fixed external input.
  \item An embodiment-agnostic adaptation formulation and a controlled AgiBot G2
  study in which the target robot is absent from the VLA foundation models'
  pretraining data. The study separates generic simulation-only policy
  adaptation from the effect of scene-native asset provenance, and measures
  zero-real-demonstration success, sim--real predictiveness, and
  real-demonstration saving on the primary long-horizon benchmark.
\end{itemize}

\section{Related Work}
\label{sec:related-work}

\subsection{Robot Simulation and Gaussian Rendering}
3DGS provides real-time novel-view rendering with explicit Gaussian
primitives~\cite{kerbl2023gaussians}. SplatSim established zero-shot transfer of
RGB manipulation policies~\cite{qureshi2025splatsim}, while Re3Sim combined
photorealistic reconstruction with domain randomization and mixed simulation
assets~\cite{han2025re3sim}. GSWorld introduced a portable Gaussian-on-mesh asset
format and a closed-loop suite for training, evaluation, and DAgger data
collection~\cite{jiang2025gsworld}. Over the past year, GS-Playground scaled
batched Gaussian rendering and rigid-body physics for visual RL
~\cite{jia2026gsplayground}, and LEGS used a native 3DGS background with a mesh
foreground to fine-tune humanoid VLAs without teleoperation~\cite{kim2026legs}.
These works advance simulator scope, rendering throughput, or downstream
learning. \system{} instead presents an embodiment-agnostic adaptation pipeline
and studies how the provenance of interactive visual assets affects the value of
data used to adapt a general-purpose VLA foundation model to an unseen target
robot.

\subsection{Real2Sim Assets and Interactive Digital Twins}
Real2Render2Real reconstructs separately scanned rigid and articulated objects
and expands one human video into robot demonstrations without dynamics
simulation~\cite{yu2025real2render2real}. Mirage2Matter separately reconstructs
scene and object Gaussians, aligns generated collision meshes, and trains VLAs
with physics-grounded data~\cite{gao2026mirage2matter}. A recent scene-level
pipeline lifts multi-view semantics into 3DGS and extracts collision meshes for
a perception--planning--execution workflow with real-robot validation
~\cite{sun2026digitaltwin}. In contrast, \system{}
requires one jointly optimized scene reconstruction to supply both the static
background and movable native-Gaussian objects, and explicitly completes both
the newly exposed background and unobserved object surfaces.

Beyond rigid scene assets, recent work estimates inertial parameters through
robot interaction~\cite{pfaff2025scalablereal2sim}, builds Gaussian soft-body
twins and studies paired sim--real policy correlation on deformable and rigid
tasks~\cite{zhang2025softbodyeval}, maintains online
prediction--correction twins for rigid bodies and deformable linear
objects~\cite{cai2026gausstwin}, and recovers articulated Gaussian assets from
static and dynamic RGB-D sequences~\cite{mishra2026artitwinsplat}. SplatCtrl instead couples online
Gaussian reconstruction to reactive collision avoidance~\cite{jain2026splatctrl}.
These contributions are complementary: \system{} assumes rigid objects and
treats physical parameters as fixed inputs. Unlike articulation discovery from
observed object motion, it extracts freely movable rigid objects from one joint
scene reconstruction and measures their value for VLA adaptation.

\subsection{Simulation Data for VLA Adaptation}
VLAs such as \piHalf{}~\cite{black2025pi05} and
SmolVLA~\cite{shukor2025smolvla} support post-training on task-specific robot
data. Real2Render2Real, Mirage2Matter, and LEGS show that reconstructed or
hybrid-rendered data can train policies that transfer to hardware; their
evaluations include fixed-budget, data-scale, and collection-cost comparisons.
We ask a different question: after holding trajectories, physics, cameras,
visual-calibration effort, and optimization fixed, how does asset provenance
shift the complete mixed sim--real learning curve when adapting a foundation
model to an unseen embodiment? This design estimates the real-demonstration
equivalent of scene-native simulation data rather than only comparing one or two
data budgets.

\section{\system: Scene-Native Interactive Gaussian Adaptation}
\label{sec:method}

\subsection{Problem Formulation}
\label{sec:problem}

The input is a calibrated multi-view capture
$\mathcal{V}=\{(I_j,\mathbf{K}_j,\mathbf{T}_j)\}_{j=1}^{J}$ of a target
workspace, where $I_j$ is an RGB image and $(\mathbf{K}_j,\mathbf{T}_j)$ are
camera intrinsics and extrinsics. The reconstructed scene is a set of Gaussian
primitives
\begin{equation}
  \mathcal{G}_{\mathrm{scene}}
  = \mathcal{G}_{\mathrm{bg}} \cup \mathcal{G}_{\mathrm{robot}}
  \cup \bigcup_{k=1}^{K}\mathcal{G}_{\mathrm{obj}}^{k},
  \label{eq:scene-decomposition}
\end{equation}
where each Gaussian
$g_i=(\boldsymbol{\mu}_i,\boldsymbol{\Sigma}_i,\alpha_i,
\mathbf{c}_i,z_i,r_i)$ contains a mean, covariance, opacity, appearance
coefficients, instance label $z_i$, and provenance label
$r_i\in\{\mathrm{observed},\mathrm{synthesized}\}$. \system{} infers $z_i$ and
constructs an interactive asset
\begin{equation}
  \mathcal{A}_k =
  (\mathcal{G}_{\mathrm{obj}}^{k},\mathcal{M}_k,\mathcal{P}_k)
  \label{eq:asset}
\end{equation}
for each object, where $\mathcal{M}_k$ is a watertight collision mesh and
$\mathcal{P}_k$ contains physical parameters supplied by a fixed external
estimator. The scene and assets are rigid, metrically scaled, and expressed in a
common simulator frame. We assume calibrated cameras and do not model dynamic
relighting, transparent objects, or deformable bodies.

At time $t$, an embodiment-agnostic VLA policy receives rendered camera
observations $\mathbf{o}_t$,
target-robot proprioception $\mathbf{s}_t$, and a language instruction $\ell$. It predicts
an $H$-step action chunk
\begin{equation}
  \mathbf{a}_{t:t+H-1}
  = \pi_{\theta}(\mathbf{o}_t,\mathbf{s}_t,\ell).
  \label{eq:vla-policy}
\end{equation}
A frozen target-robot execution stack $\mathcal{E}_{\mathrm{robot}}$ maps the first
$k\leq H$ actions to executable commands, applies the same low-level controller
and safety checks in every condition, and then acquires a new observation. The
special case $k=H$ corresponds to chunk-level replanning. \todo{Specify the
action parameterization, $H$, $k$, policy frequency, low-level control
frequency, and safety checks used in the actual system.}

Let $\mathcal{D}_{\mathrm{sim}}^m$ be the simulation-generated trajectories
under asset condition $m$, and let $\mathcal{D}_{\mathrm{real}}(n)$ contain $n$
real target-robot demonstrations. For a target real-world success rate $\tau$,
$N_{\mathrm{real}}^{m}(\tau)$ denotes the minimum $n$ for which post-training on
$\mathcal{D}_{\mathrm{sim}}^m\cup\mathcal{D}_{\mathrm{real}}(n)$ reaches
$\tau$. We measure the real-demonstration value of Scan2VLA data as
\begin{equation}
  \mathrm{Saving}(\tau) = 1 -
  \frac{N_{\mathrm{real}}^{\mathrm{Scan2VLA}}(\tau)}
       {N_{\mathrm{real}}^{\mathrm{RealOnly}}(\tau)}.
  \label{eq:saving}
\end{equation}
This definition separates real-demonstration efficiency from peak policy score
and from simulator-construction cost. We report capture, reconstruction,
calibration, and compute costs separately rather than treating zero real
demonstrations as zero real-world information.

\subsection{System Overview}
Figure~\ref{fig:pipeline} organizes \system{} along the data flow. First, scene
construction jointly optimizes one metric 3DGS and lifts multi-view instance
evidence to Gaussian labels. Second, asset construction applies separate object
and background completion, records which primitives are observed or synthesized,
and extracts collision geometry aligned to each movable Gaussian subset. Third,
the simulator transforms visual and collision representations from the same
object pose, generates synchronized image--language--state--action trajectories,
  and post-trains the foundation-model VLA. The reconstruction and completion modules operate only
during digital-twin construction; the VLA is the only learned component updated
during policy post-training. Camera calibration, physical parameters, the G2
action adapter, low-level controller, and safety stack remain fixed across all
  asset conditions; the target-robot interface adapter and execution stack are
  fixed deployment components rather than learned contributions.

\begin{figure*}[t]
  \centering
  \fbox{\parbox[c][3.2cm][c]{0.96\textwidth}{\centering
    Method pipeline with five blocks: (1) calibrated multi-view capture and
    native 3DGS; (2) multi-view instance evidence and Gaussian label fusion;
    (3) object/background completion with observed/synthesized provenance;
    (4) watertight collision geometry and physics binding; (5) synchronized
    trajectory generation, foundation-model VLA post-training, and frozen G2
    execution. Mark
    optimized, pretrained/fixed, and post-trained modules with distinct styles.}}
  \caption{Overview of \system{}. A single jointly optimized scene is
  decomposed into independently movable Gaussian assets. Geometry completion
  enables physical interaction, while the original observed Gaussians preserve
  capture-conditioned appearance.}
  \label{fig:pipeline}
\end{figure*}

\subsection{Native Scene Reconstruction}
Independent object scans create separate coordinate and capture domains. To
avoid this split, \system{} captures the target workspace and all interactive
objects together under the deployment appearance, then optimizes a single 3DGS
$\mathcal{G}_{\mathrm{scene}}$ from $\mathcal{V}$. Camera poses define a shared
scene frame; \todo{metric-scale source, such as a calibrated robot geometry or
fiducial} fixes scale and gravity, and a rigid transform aligns this frame with
the simulator and G2 base. Because every object Gaussian is optimized jointly
with its surrounding table and background, extraction does not require texture
transfer or a second photometric alignment stage. \todo{State whether G2 is
present during capture; report camera model, number of views, capture time,
3DGS objective and hyperparameters, optimization time, and frame-alignment
error.}

\subsection{Instance-Level Gaussian Decomposition}
\label{sec:decomposition}
Moving an object requires instance labels on Gaussian primitives rather than
view-specific pixels. Let $M_j^k$ be the 2D mask of instance $k$ in view $j$,
$\Pi_j$ the camera projection, $v_{ij}$ a visibility indicator, and $w_{ij}$ a
  confidence weight. \system{} aggregates the support for assigning Gaussian $i$ to
instance $k$ as
\begin{equation}
  p_i(k)=
  \frac{\sum_j v_{ij}w_{ij}
  \mathbb{I}[\Pi_j(\boldsymbol{\mu}_i)\in M_j^k]}
  {\sum_j v_{ij}w_{ij}+\epsilon}.
  \label{eq:instance-vote}
\end{equation}
It assigns $z_i=\arg\max_k p_i(k)$ when the maximum score exceeds threshold
$\eta$ and otherwise retains the primitive as background or uncertain. A
spatial-connectivity and boundary-cleanup stage removes isolated labels and
table/background leakage before estimating a canonical rigid frame for each
object. The output is one movable subset $\mathcal{G}_{\mathrm{obj}}^k$ per
object, a robot subset, a static background subset, and per-Gaussian confidence
for later failure analysis. \todo{Specify the 2D mask model or prompts,
visibility test, confidence weight, $\eta$, connectivity rule, canonical-frame
estimator, and whether any manual correction is permitted.}

\subsection{Object and Background Completion}
Extraction alone is insufficient: removing an object reveals a hole in the
background, while rotating it reveals its previously occluded underside and
  rear surfaces. \system{} therefore applies two completion operators with different
supports,
\begin{equation}
  \widehat{\mathcal{G}}_{\mathrm{bg}}
  =\mathcal{G}_{\mathrm{bg}}^{\mathrm{obs}}
  \cup\mathcal{G}_{\mathrm{bg}}^{\mathrm{syn}},\qquad
  \widehat{\mathcal{G}}_{\mathrm{obj}}^k
  =\mathcal{G}_{\mathrm{obj}}^{k,\mathrm{obs}}
  \cup\mathcal{G}_{\mathrm{obj}}^{k,\mathrm{syn}}.
  \label{eq:completion}
\end{equation}
Background completion is restricted to rays formerly occluded by the extracted
object and must connect continuously to the observed support surface. Object
completion operates in the canonical object frame and must produce appearance
and geometry that remain valid under rigid motion. Newly created primitives are
tagged as synthesized through $r_i$, allowing visual and geometric metrics to be
reported separately on observed and completed regions rather than hiding
hallucinated surfaces in an aggregate score. \todo{Specify the actual image/depth
completion models, supervision, optimization losses, multi-view consistency
mechanism, and rejection criteria for failed completion.}

\subsection{Collision Geometry and Physics Binding}
Photorealistic Gaussians do not by themselves define stable contact geometry.
For each completed object, \system{} extracts a surface from
$\widehat{\mathcal{G}}_{\mathrm{obj}}^k$, removes low-confidence and disconnected
components, closes remaining holes, and simplifies the resulting watertight mesh
for collision decomposition. The visual subset and collision mesh share the
same canonical frame, so a simulator pose updates both without an additional
alignment. We measure this binding by the bidirectional surface distance and by
collision false-positive/false-negative rates at sampled contact queries.
Physical parameters $\mathcal{P}_k$ are supplied by our separate prior method
\todo{cite the physical-property paper} and remain identical across visual-asset
conditions; this paper does not claim system identification. \todo{Report the
surface extractor, opacity/confidence thresholds, component filter, hole-closing
method, target mesh complexity, convex-decomposition settings, and asset import
format.}

\subsection{Dynamic Gaussian Transformation}
For a rigid object pose $(\mathbf{R}_k^t,\mathbf{t}_k^t)$, its Gaussian means
and covariances evolve as
\begin{equation}
  \boldsymbol{\mu}_i^t = \mathbf{R}_k^t\boldsymbol{\mu}_i^0
  + \mathbf{t}_k^t, \quad
  \boldsymbol{\Sigma}_i^t = \mathbf{R}_k^t
  \boldsymbol{\Sigma}_i^0(\mathbf{R}_k^t)^\top.
  \label{eq:rigid-transform}
\end{equation}
This transformation preserves captured radiance but does not perform
physically correct relighting; shadows and highlights may remain baked into the
representation.

\subsection{General-Purpose VLA Adaptation and Post-Training}
The simulator generates expert state--action trajectories using \todo{motion
planning, simulation teleoperation, scripted skills, or the actual combination}.
At every policy step it records
$d_t=(\mathbf{o}_t,\mathbf{s}_t,\ell,
\mathbf{a}_{t:t+H-1})$ with synchronized camera images, target-robot state,
instruction,
and action chunk. A trajectory is retained only under a prespecified validity
rule; the same accepted state--action sequence is rendered with \system{},
Independent-GS, GS-background/mesh-object, and mesh/PBR assets. Thus the visual
asset source changes while task states, actions, collision geometry, physical
parameters, initial-state distribution, and frame count remain fixed.

For each foundation-model backbone, post-training uses the same action representation,
normalization, image and language preprocessing, augmentation, trainable
parameters, optimizer steps, and checkpoint-selection rule. At deployment, the
policy consumes real target-robot observations and the frozen execution stack applies
Eq.~\eqref{eq:vla-policy} using the same action semantics as simulation.
\todo{Report trajectory counts and duration, successful/failed trajectory
ratio, train/validation split by scene and trajectory, optimizer and learning
rate, batch size, training steps, model-selection split, and inference
preprocessing.}

Before applying any experimental treatment, verify that G2 is absent from each
starting foundation model's pretraining and adaptation data. Record the exact
checkpoint hash, pretraining-data manifest and cutoff, and the operational
definition of ``G2-specific.'' Report all G2 information used outside the VLA,
including calibration, URDF/geometry, action adapters, and low-level controller
training. Base-zero-shot receives no optimization; it shares the same checkpoint
architecture and frozen execution interface as the post-trained conditions. All
post-trained conditions use identical trainable parameters and optimizer budgets
so that comparisons isolate training data rather than architecture changes.

\section{Experiments}
\label{sec:experiments}

In this section, we validate the \system{} framework. To analyze its
performance, we design a set of experiments to answer the following questions:
(i)~\textbf{RQ1: Asset validity.} Can automatic scene decomposition and
bidirectional completion produce
accurate, collision-ready object and background assets?
(ii)~\textbf{RQ2: Asset provenance and predictiveness.} Does preserving a
movable object's scene-native Gaussian appearance improve paired sim--real
visual consistency over independently reconstructed Gaussian and mesh-rendered
assets, and does simulation performance predict real performance across paired
validation groups? (iii)~\textbf{RQ3: Policy-data value.} Can simulation-generated
\system{} trajectories adapt a general-purpose VLA foundation model to an unseen
AgiBot G2 with zero real demonstrations,
and how many real demonstrations does it save at a fixed real-world success rate
across \piHalf{} and SmolVLA? Finally, (iv)~\textbf{RQ4: Scope and failure.}
Does the effect persist from a short-horizon bottle-to-basket task to whole-body,
long-horizon behavior, which components and scene variations drive it, and where
does the scene-native assumption break down?

RQ1 is tested by decomposition, completion, geometry, and collision metrics;
RQ2 by paired rendering, paired policy rollouts, and matched asset-source
controls; RQ3 by zero-real-data
comparisons and real-data learning curves; and RQ4 by stage-wise/end-to-end task
results, pose and scene stratification, efficiency measurements, and failure
counts. Each main figure and table is assigned to one of these questions.

\subsection{Experimental Setup}
\subsubsection{Robot and policy interface.}
The formulation is independent of the target embodiment; our real-robot
evaluation uses the AgiBot G2. All VLA post-training runs on an NVIDIA
RTX A6000 GPU; deployment inference runs on an NVIDIA RTX 5060 connected to the
G2 execution system. We report the GPU software stack and whether the inference
host is onboard or external, together with its communication interface. We also
report the robot's active degrees of freedom,
camera views, proprioceptive state, end effector, VLA action representation,
action-adapter mapping, and the relationship between high-level policy updates
and the low-level locomotion/manipulation controller. We also report policy and
control frequencies, mean and P95 observation-to-action latency, action-chunk
length, replanning interval, and all safety-stop conditions. A checkpoint and
training-manifest audit establishes that G2 is absent from the foundation models'
pretraining and adaptation data; this claim explicitly excludes the frozen
G2-specific execution stack. \todo{Fill the remaining CPU, memory, CUDA/software,
G2 action/control, latency, communication, and checkpoint-provenance values.}

\subsubsection{Tasks and success criteria.}
The short-horizon diagnostic task requires G2 to grasp one beverage bottle from
the table and place it inside the picnic basket. It succeeds only if the correct
bottle is stably lifted, released within the basket target volume, and remains
there for \todo{hold time} without intervention, timeout, prohibited collision,
drop, or safety abort. \todo{Specify the bottle, instruction, initial pose
distribution, basket target volume, lift threshold, episode horizon, and pose
tolerance.} The primary
whole-body task begins with snack items distributed on a cluttered table. G2
must identify and grasp every target item, place the items one by one inside a
picnic basket, then use whole-body motion to lift and carry the loaded basket and
place it fully inside a car trunk within \todo{episode time limit}. We report
cumulative success for item selection/grasping, packing all required items,
whole-body basket lift/carry, and trunk placement. An episode counts as
end-to-end success only if every required item remains inside the basket and the
loaded basket reaches the prespecified stable pose inside the trunk without
human intervention, timeout, prohibited collision, item drop, basket drop, or
safety abort. \todo{Specify the snack set, item count, tabletop pose distribution,
basket and trunk target volumes, stable-placement tolerance, and time limit.}

End-to-end success is the primary endpoint, but it is not the sole long-horizon
metric. For $K$ required items, we report item completion
$C_{\mathrm{item}}=K^{-1}\sum_{k=1}^{K}\mathbb{1}[o_k\text{ is in the basket}]$
both before basket lift and after trunk placement. Table~\ref{tab:task-metrics}
defines the remaining process, quality, efficiency, and safety measurements.

\begin{table*}[t]
  \caption{Long-horizon task measurements. Binary stage completion is reported
  with successes/trials; process metrics are summarized per episode with units
  and uncertainty.}
  \label{tab:task-metrics}
  \centering
  \small
  \begin{tabular}{p{0.16\textwidth}p{0.33\textwidth}p{0.43\textwidth}}
    \toprule
    Phase & Completion event & Process and quality measurements \\
    \midrule
    Item selection/grasp & Correct snack lifted clear of the table &
      Correct-object rate; grasp attempts per item; regrasp, wrong-object, and
      drop counts; time to stable lift \\
    Basket packing & All required snacks lie inside the basket target volume &
      $C_{\mathrm{item}}$; placement error and containment margin; sequence
      completion; per-item time; collisions with basket or neighboring items \\
    Whole-body lift/carry & Loaded basket reaches the trunk approach region &
      Basket lift height; peak/RMS roll and pitch; retained-item fraction;
      carry distance and duration; tracking error; balance/safety events \\
    Trunk placement & Basket is fully contained and stable in the trunk &
      Final position/orientation error; containment margin; stable hold time;
      final retained-item fraction; contact/collision count \\
    Episode/system & All phase events hold without a disqualifying event &
      Total and per-phase duration; completed stages; first-failure phase and
      cause; action chunks/replans; mean/P95 RTX 5060 inference and E2E latency \\
    \bottomrule
  \end{tabular}
\end{table*}

\subsubsection{Paired sim--real validation groups.}
The three validation groups assigned to each model are independent paired
evaluation groups, not training seeds. Three tasks, three scenes, and two VLA
backbones define 18 task--scene--backbone models; three groups per model yield
54 real validation units. Each group contains four real initial-state seeds, so
each model contributes $3\times4=12$ real episodes and the real evaluation pool
contains $18\times12=216$ episodes. Sim--real correlation does not require
simulating the entire real pool. Before observing simulation outcomes, we select
a stratified subset $\mathcal{P}$ of the 216 episodes, preserving coverage of all
tasks, scenes, backbones, and validation groups, and replay only those initial
states in simulation. The same selected seed is used unchanged in both domains.
The paired simulation cost is therefore $N_{\mathrm{pair}}=|\mathcal{P}|$, not
an additional fixed multiple of 216. \todo{Prespecify $N_{\mathrm{pair}}$, its
allocation across the 18 models and 54 groups, and the selection seed.}

For model $m\in\{1,\ldots,18\}$, let $\mathcal{P}_m$ be its selected paired
episodes and $n_m=|\mathcal{P}_m|$. We compute
\begin{equation}
  \widehat{p}^{\mathrm{real,pair}}_m
    =\frac{1}{n_m}\sum_{i\in\mathcal{P}_m}y_i^{\mathrm{real}},
  \qquad
  \widehat{p}^{\mathrm{sim,pair}}_m
    =\frac{1}{n_m}\sum_{i\in\mathcal{P}_m}y_i^{\mathrm{sim}}.
  \label{eq:sim-real-rates}
\end{equation}
The same task definition, initial-state distribution, success predicate, episode
horizon, and frozen execution interface are used in both domains. Training seeds,
if repeated, form a separate experimental factor and are never called validation
groups.

\subsubsection{Protocol and data splits.}
Across asset conditions, we freeze the G2 controller and action adapter and
match expert trajectories, task decomposition, resets, physical parameters,
camera poses, and trajectory-filtering rules. Training, validation, and test
sets are split by trajectory and scene to prevent adjacent-frame or trajectory
leakage. Initial object poses, distractors, and method order are randomized from
prespecified distributions. \todo{Specify whether object, basket, and trunk
poses are measured by calibrated tags, motion capture, robot perception, or
blinded video annotation, and report measurement resolution and missing-data
handling.} We report \todo{number of twins, scenes, objects,
simulation trajectories, seeds, and real trials} as successes over total trials
for every method and task. Human intervention, timeout, collision, safety
rejection, and emergency stop count as failures; reset operators are blind to
the next method whenever feasible. Compute, capture time, reconstruction time,
RTX A6000 post-training time, RTX 5060 deployment latency, and wall-clock
real-robot evaluation time are reported separately.

\subsection{Baselines and Controlled Comparisons}
\begin{table*}[t]
  \caption{Controlled adaptation conditions. All post-trained conditions share
  starting VLA weights, trajectories, action labels, trainable parameters,
  optimizer budget, collision geometry, physical parameters, camera protocol,
  initial-state distribution, and frozen target-robot execution stack.}
  \label{tab:conditions}
  \centering
  \small
  \resizebox{\textwidth}{!}{%
  \begin{tabular}{llll}
    \toprule
    Condition & Background appearance & Interactive-object appearance
      & Purpose \\
    \midrule
    Base zero-shot & Real & Real & Unadapted G2 checkpoint \\
    Real-only & Real & Real & Real-data learning curve \\
    Scan2VLA-only & Native 3DGS & In-situ native 3DGS & Zero-real-demo G2 adaptation \\
    Independent-GS-only & Native 3DGS & Separately reconstructed GS & Inserted-asset control \\
    GS-BG + Mesh-Obj-only & Native 3DGS & Mesh/PBR & Foreground-representation control \\
    Mesh/PBR-only & Mesh/PBR & Mesh/PBR & Conventional zero-real-demo control \\
    Scan2VLA + Real & Native 3DGS & In-situ native 3DGS & Real-data saving \\
    Independent-GS + Real & Native 3DGS & Separately reconstructed GS & Asset-source learning curve \\
    GS-BG + Mesh-Obj + Real & Native 3DGS & Mesh/PBR & Object-appearance learning curve \\
    Mesh/PBR + Real & Mesh/PBR & Mesh/PBR & Conventional learning curve \\
    Single-Scene & One native 3DGS & In-situ native 3DGS & Scene-diversity effect \\
    \bottomrule
  \end{tabular}
  }
\end{table*}

Table~\ref{tab:conditions} defines the main controls. We use calibrated, competitive
PBR assets rather than a deliberately weak renderer. The Independent-GS
baseline uses the same object and target scene but reconstructs the object
separately before insertion. It receives the same collision mesh, alignment
procedure, color-calibration budget, camera model, and manual intervention
budget as \system{}, matching the independent scene/object route used by systems
such as GSWorld and Mirage2Matter as closely as implementation permits. It is a
controlled independent-asset-provenance baseline, not a claimed end-to-end
reimplementation of either system. Every zero-real-data condition receives the
same simulated trajectory count and post-training budget. \todo{For each
baseline, state whether it uses official code/assets, released weights, or our
reimplementation, and report the shared hyperparameter-tuning budget.}

\subsection{Asset Reconstruction Quality}
This experiment addresses RQ1 before policy training. Instance mask IoU and
boundary F-score measure label accuracy; Gaussian leakage is the fraction of
object-labeled primitives outside the ground-truth object support. Background
residual error is computed only in the region revealed after object removal.
Object geometry is evaluated with bidirectional Chamfer distance, volume error,
and watertight reconstruction rate, while collision false-positive and
false-negative rates are measured over sampled contact queries. Metrics are
stratified by $r_i$ so observed and synthesized surfaces are not conflated.
Automatic masks are compared with ground-truth masks to separate decomposition
from completion error. \todo{Specify ground-truth acquisition, units, sample
counts, aggregation level, and uncertainty for every metric.}

\begin{figure}[t]
  \centering
  \fbox{\parbox[c][3.0cm][c]{0.94\columnwidth}{\centering
    Qualitative grid: real image, joint native 3DGS, instance decomposition,
    completed background, completed object, collision mesh, and render after a
    large object displacement.}}
  \caption{Qualitative asset construction. Highlight boundary leakage,
  background ghosts, synthesized object regions, and collision alignment rather
  than showing only successful examples.}
  \label{fig:assets}
\end{figure}

\subsection{Asset-Source Visual and Policy Consistency}
RQ2 uses paired real and simulated observations at the same object states and
camera poses. Each asset route renders the same held-out state set; LPIPS and
DINO feature distance measure full-frame perceptual discrepancy, while a
boundary-weighted error isolates the manipulated object and contact region.
Results are stratified by observed versus synthesized surface and by translation
and rotation from the captured pose. The strongest comparison is Scan2VLA versus
Independent-GS because trajectories, collision geometry, and calibration effort
are matched and only asset provenance changes. These metrics test consistency
with the capture domain, not physically correct materials or relighting.
\todo{Define the paired-state acquisition protocol, boundary band, pose bins,
number of frames, and confidence interval.}

Policy predictiveness uses the stratified paired subset in
Eq.~\eqref{eq:sim-real-rates}. The main paper reports a compact scatter plot and
Spearman's rank correlation across the 18 model-level simulation/real pairs;
Pearson correlation and mean absolute calibration error are secondary summaries.
Model-cluster bootstrap intervals preserve the three validation groups nested
within each model. At episode level, a $2\times2$ success/failure table over the
$N_{\mathrm{pair}}$ matched seeds reports both-success, sim-only-success,
real-only-success, and both-failure. The complete 54-cell ledger, subset manifest,
calibration model, and task/backbone-stratified results are supplementary.

\subsection{Real-Demonstration Efficiency}
First compare the unadapted Base-zero-shot checkpoint with Scan2VLA-only,
Independent-GS-only, GS-BG + Mesh-Obj-only, and Mesh/PBR-only at
$N_{\mathrm{real}}=0$. A supported gain for every post-trained condition over
Base-zero-shot is evidence that simulation-generated trajectories can adapt the
  G2 policy. A further Scan2VLA-only gain over the matched zero-real-data asset
controls is needed to attribute the improvement to scene-native provenance.
Report exact success counts, trial counts, and confidence intervals; do not
summarize this result only with qualitative terms such as ``strong'' or
``remarkable.'' ``Zero real demonstrations'' refers to policy post-training and
does not exclude the real-scene capture and calibration used to construct the
simulator.

For each backbone, train Real-only and Scan2VLA + Real policies with
\begin{equation}
  N_{\mathrm{real}} \in \{0,5,10,25,50\}, \qquad
  N_{\mathrm{sim}} = 500 \text{ per task}.
  \label{eq:data-budgets}
\end{equation}
\todo{Confirm the simulated and real demonstration budgets before running the
full training grid.}
Use learning-curve interpolation and bootstrap confidence intervals to estimate
Eq.~\eqref{eq:saving} at prespecified success thresholds. Compare treatment
effects within each backbone rather than interpreting absolute differences
between \piHalf{} and SmolVLA.

\begin{figure}[t]
  \centering
  \fbox{\parbox[c][3.2cm][c]{0.94\columnwidth}{\centering
    Main plot: real-robot success rate versus number of real demonstrations.
    Mark Base-zero-shot and all four simulation-only conditions at zero real
    demonstrations; show
    Real-only, Scan2VLA + Real, Independent-GS + Real, and Mesh/PBR + Real
    with 95\% confidence bands; mark the demonstration saving at threshold
    $\tau$.}}
  \caption{Primary result: real-data learning curves. The horizontal gap at a
  fixed real-world success rate defines the real demonstrations replaced by
  Scan2VLA data.}
  \label{fig:data-efficiency}
\end{figure}

\begin{table}[t]
  \caption{Main real-robot results. Report success with 95\% confidence
  intervals and the bootstrapped real-demonstration saving.}
  \label{tab:main-results}
  \centering
  \small
  \resizebox{\columnwidth}{!}{%
  \begin{tabular}{lccc}
    \toprule
    Training data & $N_{\mathrm{real}}$ & Success $\uparrow$ & Saving $\uparrow$ \\
    \midrule
    Base zero-shot & 0 & -- & -- \\
    Scan2VLA-only & 0 & -- & -- \\
    Independent-GS-only & 0 & -- & -- \\
    GS-BG + Mesh-Obj-only & 0 & -- & -- \\
    Mesh/PBR-only & 0 & -- & -- \\
    Real-only & -- & -- & -- \\
    Mesh/PBR + Real & -- & -- & -- \\
    GS-BG + Mesh-Obj + Real & -- & -- & -- \\
    Independent-GS + Real & -- & -- & -- \\
    Scan2VLA + Real & -- & -- & -- \\
    \bottomrule
  \end{tabular}
  }
\end{table}

\subsection{Long-Horizon G2 Results}
RQ4 tests whether an asset-source effect observed in bottle-to-basket pick-and-place
survives the compounding perception, manipulation, locomotion, and control
requirements of the picnic task. Table~\ref{tab:task-scale} reports successes
over trials at every stage and uses end-to-end success as the primary long-horizon
metric. It is interpreted together with Table~\ref{tab:task-metrics}: each method
also reports item completion, grasp attempts and drops, basket-pose stability,
item retention during carry, final trunk pose error, phase durations, safety
events, and the first failed phase. This prevents equal binary outcomes from
hiding materially different execution quality or failure modes.

\begin{table*}[t]
  \caption{Zero-real-demonstration G2 results. Every entry is reported as
  successes/trials with a 95\% binomial confidence interval. Packing, whole-body
  lift/carry, and trunk placement are cumulative stages; E2E requires all stages.}
  \label{tab:task-scale}
  \centering
  \small
  \begin{tabular}{lcccccc}
    \toprule
    Training data & Bottle task & Select/Grasp & Pack & Lift/Carry & Trunk & E2E \\
    \midrule
    Base zero-shot & -- & -- & -- & -- & -- & -- \\
    Mesh/PBR-only & -- & -- & -- & -- & -- & -- \\
    GS-BG + Mesh-Obj-only & -- & -- & -- & -- & -- & -- \\
    Independent-GS-only & -- & -- & -- & -- & -- & -- \\
    Scan2VLA-only & -- & -- & -- & -- & -- & -- \\
    \bottomrule
  \end{tabular}
\end{table*}

\todo{Write the result paragraph as: purpose; exact Scan2VLA and strongest
baseline successes/trials; percentage-point difference with uncertainty; effect
on item completion, attempts/drops, basket stability, final pose error, and
duration; stage at which the gap first appears; bounded interpretation; failure
regime.}

\subsection{Ablations and Generalization}
To identify mechanism rather than only final performance, we replace automatic
labels with ground-truth masks, remove object completion, remove background
completion, and vary the collision-geometry construction while keeping data and
training fixed. The first comparison isolates decomposition error; the two
completion ablations test newly exposed visual regions; geometry variants test
contact fidelity rather than rendering. Data-scale sensitivity varies
$N_{\mathrm{sim}}$ under a fixed optimizer protocol. Pose robustness is reported
over prespecified $0$--$30^\circ$, $30$--$90^\circ$, and $>90^\circ$ rotation
bins and corresponding translation bins. Target-twin adaptation, held-out scene,
new object instance, appearance change, and language paraphrase are labeled as
separate generalization axes defined before testing; no test result is used to
choose bins or hyperparameters.

\subsection{Efficiency, Safety, and Failure Analysis}
System cost is decomposed into scene capture, 3DGS optimization, decomposition,
completion, collision processing, trajectory generation, and VLA post-training.
At deployment on the RTX 5060, we report image acquisition, preprocessing, VLA
inference, action conversion, safety validation, and end-to-end latency as mean and P95, together
with achieved policy and control frequencies. Safety reporting includes
collision aborts, command rejection, stale-state timeout, emergency stops, and
human interventions; these events remain failures in task success.

Every failed real episode is assigned one primary stage and cause: perception or
wrong-object selection, grasp, packing, basket lift/carry, trunk placement,
VLA timeout, low-level tracking, safety rejection, or hardware/communication.
The paper reports count, fraction, detectability before motion, and whether the
failure occurs on observed or synthesized Gaussian regions. This analysis tests
whether the asset-source advantage degrades under large viewpoint changes,
occlusion, specular/transparent packaging, load-induced control error, or
long-horizon compounding.

\subsection{Statistical Analysis}
The independent statistical unit is an episode, not an image frame. Initial
states are paired across asset conditions and, where reproducible, across
simulation and reality. We report per-task and aggregate successes/trials with
exact binomial 95\% confidence intervals and training-seed variation. Primary
inference uses strict long-horizon end-to-end success; item completion, drop
count, basket stability, final trunk pose error, and duration are prespecified
secondary outcomes, while the remaining process metrics are diagnostic. We
report the multiple-comparison correction used for confirmatory secondary
outcomes. Method comparisons use a mixed-effects logistic model with asset condition and real-data
budget as fixed effects and task, scene, object, and initial state as random
effects~\cite{kressgazit2024empirical}; small paired samples additionally receive
an exact paired or Fisher test as appropriate. Sim--real correlation uses the
prespecified subset of paired real/simulation episodes and preserves the 54
validation cells nested within 18 models. Intervals use model-cluster bootstrap;
the supplementary hierarchical binomial model accounts for unequal selected
subset sizes $n_m$. Correlation, calibration error, and paired-episode agreement
are reported per task and backbone before any pooled result.
Estimating
Eq.~\eqref{eq:saving} bootstraps data sampling, policy training seed, real-robot
episodes, and the complete learning-curve interpolation procedure. Statistical
significance is claimed only for prespecified comparisons with reported effect
sizes and uncertainty.

\section{Discussion and Limitations}
\label{sec:limitations}
The causal interpretation depends on both levels of control. A gain shared by
all simulation-only conditions over Base zero-shot establishes the value of
simulation-generated G2 trajectories; only an additional Scan2VLA-only gain over
Independent-GS-only and mesh-foreground controls supports the scene-native
asset-provenance hypothesis. Visual and policy gains should therefore be
reported separately rather than attributing all adaptation to decomposition.

``Zero real demonstrations'' does not mean zero real-world information: Scan2VLA
uses a target-scene capture, camera/robot calibration, robot geometry, and a
frozen G2 execution stack. The experiments adapt the foundation-model policy, not
the full robot controller. The formulation is embodiment-agnostic, but the
reported evidence covers one genuinely unseen target embodiment, AgiBot G2;
transfer to another robot requires a separate interface and execution audit.
Target-twin results measure adaptation to captured deployment scenes; held-out
objects, backgrounds, lighting, language, and scenes are separate generalization
claims and must not be inferred from target-scene success.

Native 3DGS preserves capture-conditioned appearance but does not separate
material from illumination. Large object motions may expose synthesized regions
or retain baked shadows and highlights. Transparent, specular, thin, or heavily
occluded objects remain difficult; completion can create plausible but incorrect
geometry, and segmentation errors propagate to collision assets. The current
formulation assumes rigid objects and fixed externally estimated physical
parameters. Long-horizon G2 performance also depends on the frozen action
adapter, low-level controller, calibration, and safety stack, so conclusions do
not automatically transfer to another embodiment or controller.

\section{Conclusion}
\label{sec:conclusion}
We presented \system, an embodiment-agnostic pipeline that converts a jointly
reconstructed native 3DGS scene into an interactive Gaussian digital twin with
movable object assets, completed background/object regions, and aligned collision
geometry. On AgiBot G2, which is absent from the evaluated foundation models'
pretraining data, the controlled study measures simulation-only adaptation,
scene-native asset provenance, and sim--real predictiveness across short- and
long-horizon tasks. \todo{State the exact Scan2VLA-only successes/trials versus
the strongest matched baseline and the real-demonstration saving with confidence
intervals.} The evidence \todo{state only the bounded conclusion supported for
the evaluated rigid objects, captured scenes, unseen G2 embodiment, VLA
backbones, and frozen execution stack}.

\begin{thebibliography}{17}

\bibitem{kerbl2023gaussians}
B. Kerbl, G. Kopanas, T. Leimkuehler, and G. Drettakis,
``3D Gaussian splatting for real-time radiance field rendering,''
\emph{ACM Trans. Graph.}, vol. 42, no. 4, 2023.

\bibitem{qureshi2025splatsim}
M. N. Qureshi et al., ``SplatSim: Zero-shot sim2real transfer of RGB
manipulation policies using Gaussian splatting,'' in \emph{Proc. IEEE Int.
Conf. Robot. Autom. (ICRA)}, 2025.

\bibitem{han2025re3sim}
X. Han et al., ``Re3Sim: Generating high-fidelity simulation data via
3D-photorealistic real-to-sim for robotic manipulation,'' arXiv:2502.08645,
2025.

\bibitem{jiang2025gsworld}
G. Jiang et al., ``GSWorld: Closed-loop photo-realistic simulation suite for
robotic manipulation,'' arXiv:2510.20813, 2025.

\bibitem{yu2025real2render2real}
J. Yu et al., ``Real2Render2Real: Scaling robot data without dynamics
simulation or robot hardware,'' arXiv:2505.09601, 2025.

\bibitem{pfaff2025scalablereal2sim}
N. Pfaff et al., ``Scalable Real2Sim: Physics-aware asset generation via
robotic pick-and-place setups,'' arXiv:2503.00370, 2025.

\bibitem{zhang2025softbodyeval}
K. Zhang et al., ``Real-to-sim robot policy evaluation with Gaussian splatting
simulation of soft-body interactions,'' arXiv:2511.04665, 2025.

\bibitem{sun2026digitaltwin}
Z. Sun et al., ``A high-fidelity digital twin for robotic manipulation based
on 3D Gaussian Splatting,'' arXiv:2601.03200, 2026.

\bibitem{gao2026mirage2matter}
Z. Gao et al., ``Mirage2Matter: A physically grounded Gaussian world model
from video,'' arXiv:2602.00096, 2026.

\bibitem{cai2026gausstwin}
Y. Cai et al., ``GaussTwin: Unified simulation and correction with Gaussian
splatting for robotic digital twins,'' arXiv:2603.05108, 2026.

\bibitem{jia2026gsplayground}
Y. Jia et al., ``GS-Playground: A high-throughput photorealistic simulator for
vision-informed robot learning,'' arXiv:2604.25459, 2026.

\bibitem{kim2026legs}
H. Kim et al., ``LEGS: Fine-tuning teleop-free VLAs for humanoid
loco-manipulation in an embodied Gaussian splatting world,'' arXiv:2606.01458,
2026.

\bibitem{mishra2026artitwinsplat}
P. Mishra et al., ``ArtiTwinSplat: Interactable digital twin reconstruction via
Gaussian splatting from RGB-D videos,'' arXiv:2606.24628, 2026.

\bibitem{jain2026splatctrl}
S. Jain and H. J. Choi, ``SplatCtrl: Perception-action coupling via Gaussian
scene representations and reactive robot control,'' arXiv:2607.08948, 2026.

\bibitem{black2025pi05}
K. Black et al., ``$\pi_{0.5}$: A vision-language-action model with open-world
generalization,'' arXiv:2504.16054, 2025.

\bibitem{shukor2025smolvla}
M. Shukor et al., ``SmolVLA: A vision-language-action model for affordable and
efficient robotics,'' arXiv:2506.01844, 2025.

\bibitem{kressgazit2024empirical}
H. Kress-Gazit et al., ``Robot learning as an empirical science: Best practices
for policy evaluation,'' arXiv:2409.09491, 2024.

\end{thebibliography}

% Before submission, replace provisional arXiv entries with archival citations,
% remove all TODOs, inspect anonymous PDF metadata, and verify the current ICRA
% template and eight-page total limit in PaperPlaza.

\end{document}
